% !TEX root = ../../main.tex

\section{Cierre}

El trabajo muestra una vía concreta para ampliar \texttt{ApplicativeDo} sin abandonar el compilador real ni la sintaxis habitual de \textsf{Haskell}. Bajo una hipótesis explícita de conmutatividad, el compilador puede evaluar todos los reordenamientos válidos, incluido el orden escrito por el programador; por definición, todos preservan las dependencias locales de variables.

El aporte principal no es solo el algoritmo de generación de reordenamientos válidos, sino su integración experimental en \textsf{GHC} con trazabilidad: módulos agregados, parches, corpus, scripts y logs reproducibles. Esto permite estudiar la idea de manera concreta, observar sus beneficios en casos controlados y delimitar sus riesgos.

En síntesis, la memoria entrega evidencia experimental de que el espacio de planes de \texttt{ApplicativeDo} puede ampliarse para mónadas declaradas como conmutativas, obteniendo mejoras de costo sin modificar la semántica observable en los casos sintéticos evaluados.
